\documentclass[a4paper,11pt]{article}

\usepackage{latexsym}
\usepackage[empty]{fullpage}
\usepackage{titlesec}
\usepackage{marvosym}
\usepackage{verbatim}
\usepackage{enumitem}
\usepackage[pdftex]{hyperref}
\usepackage{fancyhdr}
\usepackage{graphicx}
\usepackage{float}
\usepackage{wrapfig}

\pagestyle{fancy}
\fancyhf{}
\fancyfoot{}
\renewcommand{\headrulewidth}{0pt}
\renewcommand{\footrulewidth}{0pt}

\addtolength{\oddsidemargin}{-0.530in}
\addtolength{\evensidemargin}{-0.375in}
\addtolength{\textwidth}{1in}
\addtolength{\topmargin}{-.45in}
\addtolength{\textheight}{1in}

\urlstyle{rm}
\hypersetup{colorlinks=false,pdfborder={0 0 0}}
\raggedbottom
\raggedright
\setlength{\tabcolsep}{0in}

\titleformat{\section}{
  \vspace{-7pt}\scshape\raggedright\large
}{}{0em}{}[\titlerule \vspace{-5pt}]

\newcommand{\resumeItem}[2]{
  \item\small{\textbf{#1}{: #2 \vspace{-1pt}}}
}
\newcommand{\resumeItemWithoutTitle}[1]{
  \item\small{{\vspace{-2pt}}}
}
\newcommand{\resumeSubheading}[4]{
  \vspace{-1pt}\item
    \begin{tabular*}{0.97\textwidth}{l@{\extracolsep{\fill}}r}
      \textbf{#1} & #2 \\
      \textit{\small #3} & \textit{\small #4} \\
    \end{tabular*}\vspace{-5pt}
}
\newcommand{\resumeSubItem}[2]{\resumeItem{#1}{#2}\vspace{-1pt}}
\renewcommand{\labelitemii}{$\circ$}
\newcommand{\resumeSubHeadingListStart}{\begin{itemize}[leftmargin=*, itemsep=2pt, parsep=0pt]}
\newcommand{\resumeSubHeadingListEnd}{\end{itemize}}
\newcommand{\resumeItemListStart}{\begin{itemize}[leftmargin=16pt, itemsep=3pt, parsep=0pt]}
\newcommand{\resumeItemListEnd}{\end{itemize}\vspace{-5pt}}
\newcommand{\resumeSkillListStart}{\begin{itemize}[leftmargin=*, itemsep=5pt, parsep=0pt]}
\newcommand{\resumeSkillListEnd}{\end{itemize}}
\usepackage[utf8]{inputenc}
\usepackage[T2A]{fontenc}

\begin{document}

%----------HEADING-----------------
\begin{tabular*}{\textwidth}{l@{\extracolsep{\fill}}r}
  \textbf{\LARGE Александр Суриков} & Telegram: \href{https://t.me/aiex_su}{@aIex\_su} \\
  Продуктовый аналитик & \\
  Москва, Россия \;|\; Яндекс.Директ & \\
\end{tabular*}
\vspace{-5pt}

\section{Навыки}
  \resumeSkillListStart
    \resumeSubItem{Главная экспертиза}{A/B-тестирование, продуктовая стратегия, causal inference, SQL/Python, рекламные продукты}
    \resumeSubItem{Продуктовая аналитика}{дерево метрик, поиск точек роста, оценка потенциала и инкрементального эффекта, приоритизация портфеля, unit-экономика рекламных сценариев}
    \resumeSubItem{Эксперименты и каузальная оценка}{A/B-дизайн, основные и защитные метрики, CUPED, стратификация, bucketed t-test и Mann-Whitney, MDE, мощность, глобальный контроль, propensity-модели}
    \resumeSubItem{Экспертиза и влияние}{ответственность за аналитику направления, декомпозиция неопределенных задач, защита выводов перед бизнесом, ревью экспериментов, наставничество аналитиков, найм стажеров}
    \resumeSubItem{Инструменты}{SQL: YQL, ClickHouse; Python: pandas, NumPy, scipy; YT, MapReduce, Nirvana/Airflow, DataLens/Tableau, Git, Bash/Linux}
  \resumeSkillListEnd
\vspace{0pt}

\section{Опыт работы}
  \resumeSubHeadingListStart

\resumeSubheading{Яндекс}{2023 -- текущее время}
    {Продуктовый аналитик, Яндекс.Директ}{}
    \par\vspace{7pt}
    {\small\textbf{Стек:} \textit{YQL, ClickHouse, Python, pandas, scipy, A/B и causal inference, YT, Nirvana, DataLens}}
    \par\vspace{1pt}
    \resumeItemListStart
        \resumeItem{Выстроил аналитику направления Рекомендаций}{собрал дерево метрик, глобальный контроль для оценки инкрементального эффекта и стандарты A/B-дизайна. Масштаб -- сотни тысяч клиентов и десятки тысяч применений в месяц; подтвержденный эффект около 100 млн руб. в месяц}
\vspace{2pt}
        \resumeItem{Запускал подсказки по стартовым настройкам кампаний}{оценил потенциал и сопроводил запуск подсказок по недельному бюджету и CPA-настройкам. Доля новых кампаний, продолжающих работу после старта, выросла на 7\%, средняя выручка новых кампаний -- более чем на 5\%}
\vspace{2pt}
        \resumeItem{Отвечал за аналитику редизайна страницы рекомендаций}{определил метрики и A/B-дизайн с контролем каннибализации и эффекта новизны. Эксперимент показал +40\% к применениям рекомендаций по продукту и +60\% на ключевой поверхности}
\vspace{2pt}
        \resumeItem{Вел портфель рекомендаций и подсказок рекламодателям}{сопровождал 20+ A/B: формулировал гипотезы, выбирал метрики, считал MDE и доводил эксперименты до решения. Приоритизировал портфель по ожидаемому инкрементальному эффекту}
\vspace{2pt}
        \resumeItem{Стандартизировал запуск A/B}{ввел шаблон гипотез, дизайна, автометрик и критериев решения. Подготовка эксперимента сократилась с нескольких дней до часов; шаблон использован в 20+ запусках}
\vspace{2pt}
        \resumeItem{Увеличил мощность экспериментов через Propensity to Apply}{построил CatBoost-модель склонности рекламодателя применить рекомендацию и использовал ее для отбора аудитории и снижения шума в A/B. Подход повышает чувствительность экспериментов без роста трафика}
\vspace{2pt}
        \resumeItem{Усиливал аналитическую экспертизу команды}{онбордил и менторил 5 аналитиков и стажеров, ревьюил дизайн экспериментов и выводы, провел 20+ собеседований; зафиксировал процессы в гайдах}
\vspace{2pt}
        \resumeItem{Перевел мониторинг продукта в проактивный режим}{внедрил алерты на продуктовые деградации, качество данных и состояние пайплайнов. Система выявляет десятки инцидентов в год до существенного ущерба}
\vspace{2pt}
        \resumeItem{Ускорил цикл разработки рекомендаций}{перенес сборку баз из разовых YQL-запросов в продуктовые Python-пайплайны: время сборки и теста сократилось примерно с 2 часов до 30 минут}
    \resumeItemListEnd

\vspace{5pt}
\resumeSubheading{Сбер}{2022 -- 2023}
{Аналитик данных}{}
    \par\vspace{7pt}
    {\small\textbf{Стек:} \textit{SQL, Python, Hadoop}}
    \par\vspace{1pt}
\resumeItemListStart
        \resumeItem{Аналитик данных банкоматов}{построил ETL-пайплайн для логов банкоматов, нашел устройства с низким качеством распознавания QR-кодов и связал проблему с версиями прошивок}
\resumeItemListEnd
\resumeSubHeadingListEnd

\vspace{-2pt}
\section{Образование}
  \resumeSubHeadingListStart
    \resumeSubheading
      {НИУ ВШЭ}{Москва, Россия}
      {Магистратура, Высшая школа бизнеса}{2025 -- 2027}
    \vspace{5pt}
    \resumeSubheading
      {НИТУ МИСИС}{Москва, Россия}
      {Бакалавриат, Информатика и вычислительная техника}{2020 -- 2024}
  \resumeSubHeadingListEnd

\end{document}
